\section{Round-four referee closure: platform-labelled cotangent functors}
\label{sec:r3-c2}

There is no unlabelled parent law spanning physically different systems.  The
universal object is a typed variational functor acting componentwise on a
disjoint union of platforms.  This section constructs the strict dual, proves
the required continuity or exponential approximation of the mechanical
maps, derives pressure derivatives from the proved spectral or cluster
packets, and supplies explicit likelihood, diffusion, BSDE, and memory
commutation statements.

\subsection{The platform coproduct}

Let \(\mathfrak A\) index the platforms actually constructed in the series:
the periodic Sinai path platform, the A1 Hamiltonian impact platform, and the
hard-sphere Boltzmann--Grad platform.  Define
\[
 \mathfrak P
 =\bigsqcup_{\alpha\in\mathfrak A}
 \{\alpha\}\times\mathcal P_w(\Omega_\alpha),
 \qquad
 \mathfrak I(\alpha,\nu)=I_\alpha(\nu).
\]
The label has the discrete topology.  No cost is attached to the label because
no probability law randomly changes the physical platform.

\begin{theorem}[Typed contraction functor]
\label{thm:r3-c2-functor}
Let \(C_\alpha:\mathcal P_w(\Omega_\alpha)\to E_\alpha\) be continuous, or
admit an exponentially good continuous approximation under the microscopic
laws of platform \(\alpha\).  Then
\[
 I_{\alpha,C}(a)
 =\inf\{I_\alpha(\nu):C_\alpha(\nu)=a\}
\]
is the contraction rate on that platform.  A map between two different
labels may be used only when a separate exponentially good scaling theorem
between their microscopic laws has been proved.
\end{theorem}

\begin{proof}
On a fixed labelled component this is the ordinary contraction theorem, with
the extended version for exponentially good approximations.  The coproduct
topology makes every component open and closed, so no minimizing sequence can
change labels.  A cross-platform conclusion therefore requires an actual map
between the two families of laws; notation alone cannot create one.
\end{proof}

\subsection{Mechanical maps on the defect-completed Sinai platform}

The A3 parent object is the direct graph-completed collision/physical path
law, including survivor phases.  Since finite horizon makes the one-flight
roof bounded and the lattice displacement has finite range, homological
current and physical clock are bounded one-step observables.  For a fixed path
window, remove only a neighborhood of the finitely many singularity faces
which enter that window; no return-length event is discarded.

\begin{lemma}[Direct-pressure continuity ledger]
\label{lem:r3-c2-continuity}
On the periodic Sinai component, each of the following has continuous
regular-component approximants whose discrepancy has arbitrarily large cost
under the A3 direct rate:
\begin{enumerate}[label=(\roman*)]
\item finite-window collision and physical paths;
\item bounded occupation and impact marks;
\item empirical homological current and bounded roof average;
\item every bounded local path work.
\end{enumerate}
The approximation remains valid on zero-magnet-frequency survivor phases.
\end{lemma}

\begin{proof}
For a window of radius \(q\), only the finite union
\(\mathcal S_q\) of graph-completed singularity faces can destroy ordinary
continuity.  On the complement of its \(\delta\)-neighborhood, all collision
coordinates, suspension maps, finite-range displacement labels, and bounded
marks are uniformly continuous on finitely many regular components.  A3's
direct-pressure Chernoff estimate sends the cost of positive empirical mass in
that neighborhood to infinity as \(\delta\downarrow0\).  The argument is
microscopic collision-time based and applies equally to induced and survivor
branches.  Since the roof and displacement are bounded one-step observables,
no long-return truncation enters.
\end{proof}

The A1 current is continuous on its common symbolic path space.  The B2 density
and actual-contact coordinates are continuous in the energy/weighted weak
state of B2--B4.  Thus every contraction used below has a proved fixed-platform
map.

\subsection{Coercive weighted bounded-strict pressure dual}

Let \(\Omega\) be Polish and let \(W\ge1\) be continuous with compact
sublevels.  Define
\[
 C_W(\Omega)=\{F\in C(\Omega):\|F/W\|_\infty<\infty\}
\]
with the weighted bounded-strict topology \(\beta_W\).  Its dual is
\(\mathcal M_W(\Omega)\), the countably additive signed Radon measures with
finite \(W\)-moment.

\begin{theorem}[Coercive strict dual and phase compactness]
\label{thm:r3-c2-strict}
Assume the good path rate satisfies
\[
 I(\nu)\ge a\nu(W)-b
\]
for some \(a>0\).  Then for every source set
\(\|F/W\|_\infty\le a-\eta\), \(\eta>0\),
\[
 Q(F)=\sup_\nu\{\nu(F)-I(\nu)\}
\]
is finite and \(\beta_W\)-continuous on bounded subsets.  The supremum is
attained on a compact rate sublevel and
\[
 \partial_{\beta_W}Q(F)
 =\operatorname*{argmax}_\nu\{\nu(F)-I(\nu)\}.
\]
Every subgradient is a countably additive finite-weight path probability.
\end{theorem}

\begin{proof}
Weighted Riesz representation identifies the strict dual with
\(\mathcal M_W\).  The coercive inequality gives
\[
 \nu(F)-I(\nu)
 \le -\eta\nu(W)+b,
\]
so maximizing sequences have uniformly bounded \(W\)-moment and lie in one
compact rate sublevel.  This proves finiteness and attainment.  For a
bounded-strict convergent net, compact-uniform convergence controls the
integral on a compact \(W\)-sublevel and the uniform moment bound controls its
tail, proving continuity.  Fenchel equality identifies maximizers and
subgradients; testing constants and nonnegative compactly supported functions
makes every Radon subgradient a probability.
\end{proof}

On the A3 platform, the coercive inequality is the pressure-dual estimate
proved with the survivor branch.  On B2, it follows from the quadratic energy
source and collision entropy.  Thus the theorem is applied only inside an
explicit source domain, rather than inferred from goodness alone.

\subsection{The closed coboundary quotient}

On the Sinai component use the A3 dynamically H\"older weighted algebra
\(\mathscr A_\beta\) and its closed subspace
\[
 \mathscr N_\beta
 =\overline{\operatorname{span}}
 \left(
 \mathbb R\mathbf1+
 \{U-U\circ\Theta_t:t\in\mathbb Q,
 U\in\mathscr A_\beta\}
 \right).
\]

\begin{theorem}[Pressure cotangents]
\label{thm:r3-c2-cotangent}
Long-time excess pressure descends to the Hausdorff quotient
\(\mathscr A_\beta/\mathscr N_\beta\).  Two potentials have identical
integrals under every invariant finite-weight phase exactly when their
difference belongs to \(\mathscr N_\beta\).  At an exposed phase, the pressure
tangent is its countably additive path law; at coexistence the tangent set is
the compact convex set of maximizing phases from
Theorem~\ref{thm:r3-c2-strict}.
\end{theorem}

\begin{proof}
Time averages of a rational-time coboundary have a bounded endpoint remainder,
so their pressures differ only by the constant part.  The annihilator theorem
proved in A3 gives the converse.  The tangent statements follow from the
strict-dual subdifferential identity.
\end{proof}

\subsection{First and second derivatives from model packets}

Let \(F,G\) be dynamically H\"older path directions in the common complex
spectral ball of A2.  Let \(Q_R\) be the physical pressure root.

\begin{theorem}[Spectral pressure response]
\label{thm:r3-c2-response}
On every simple Sinai phase branch,
\[
 DQ_R(\Psi)F=\mathbb P_R^\Psi(F)
\]
and
\[
 D^2Q_R(\Psi)(F,G)
 =\int_{-\infty}^{\infty}
 \operatorname{Cov}_{\mathbb P_R^\Psi}
 (F,G\circ\Theta_t)\,dt.
\]
The integral is absolutely convergent and the derivatives are continuous in
\(R\).  On the B2 hard-sphere source ball, the analogous first and second
joint density--collision derivatives exist and equal the B3 mean and
covariance.  No second derivative is asserted solely from uniqueness of a
variational maximizer.
\end{theorem}

\begin{proof}
For Sinai paths, A2 gives an analytic simple eigenvalue and analytic left and
right eigenvectors for higher-block twists.  Differentiate the eigenvalue
once.  Differentiating the eigenprojector produces the reduced resolvent
series, whose exponential spectral decay sums to the bilateral Green--Kubo
integral.  Uniform spectral constants give absolute convergence and
\(R\)-continuity.  For hard spheres, B2 normal convergence and B3 microscopic
normalization give the stated derivatives.
\end{proof}

\subsection{Entropic rigidity with an operational calibration}

Let \(\mathfrak C_t\) be a monotone, cash-additive, strictly relevant,
strongly time-consistent scalar certainty equivalent generated by one
\(C^2\) utility.

\begin{theorem}[Rigidity and canonical-cocycle calibration]
\label{thm:r3-c2-rigidity}
The family is either conditional expectation or
\[
 \mathfrak C_t(X)
 =\vartheta^{-1}\log
 \mathbb E[e^{\vartheta X}\mid\mathcal F_t]
\]
with one time-independent \(\vartheta\ne0\).  The axioms alone do not choose
\(\vartheta\).  If, in addition, the likelihood-ratio cocycle is required to
be the mechanically prepared canonical cocycle generated by
\(\theta G\), and work is transported in the same cotangent bundle, then
\(\vartheta=\theta\).
\end{theorem}

\begin{proof}
Cash additivity gives the translation functional equation for the utility.
Twice differentiating shows that absolute risk aversion is constant, yielding
an affine or exponential utility.  The conditional tower makes the exponential
coefficient independent of the time split.  Under the extra canonical-cocycle
condition, applying both likelihood ratios to the generator \(G\) gives
\(e^{\vartheta G}/Ee^{\vartheta G}=e^{\theta G}/Ee^{\theta G}\); nonconstancy
of \(G\) forces \(\vartheta=\theta\).
\end{proof}

\subsection{Stable resolved filtrations and deterministic-to-diffusion likelihoods}

Let \(X^\varepsilon\) be the resolved slow coordinate in A4 and let
\[
 \mathcal G_t^\varepsilon
 =\sigma(X_s^\varepsilon:0\le s\le t)
\]
be its resolved filtration.  It is not the complete microscopic filtration.
Under a prepared Doob phase, denote the resolved transition kernel by
\(P_{s,t}^\varepsilon\).  Let \(P_{s,t}\) be the limiting diffusion kernel.

\begin{lemma}[Kernel and filtration convergence]
\label{lem:r4-c2-kernel}
For every bounded uniformly continuous \(F\) and compact resolved state set,
\[
 \sup_{x\in K}
 |P_{s,t}^\varepsilon F(x)-P_{s,t}F(x)|\to0
\]
uniformly on \(0\le s\le t\le T\).  Jointly with
\(X^\varepsilon\Rightarrow X\),
\[
 \mathbb E[F(X_T^\varepsilon)\mid\mathcal G_t^\varepsilon]
 =P_{t,T}^\varepsilon F(X_t^\varepsilon)
 \Rightarrow P_{t,T}F(X_t)
 =\mathbb E[F(X_T)\mid\mathcal G_t].
\]
The convergence holds jointly at finitely many times and in the martingale
Skorokhod topology.
\end{lemma}

\begin{proof}
A4 proves convergence of the Doob nonlinear semigroups and their linear first
derivatives on bounded cylinders.  Cylinder approximation and the energy
compact-containment estimate extend this to bounded uniformly continuous
functions uniformly on compact states.  The Markov representation gives the
two conditional expectations as transition kernels evaluated at the current
state.  Uniform kernel convergence and process convergence give joint
finite-time convergence.  The conditional processes are bounded martingales;
Aldous' criterion follows from the uniform kernel modulus, giving Skorokhod
martingale convergence.  This is a theorem about the resolved filtrations and
does not identify them with complete deterministic information.
\end{proof}

Let \(\Phi\in C_b^3(\mathbb R^d)\), \(\theta\in\mathbb R\), and define
\[
 M_t^\varepsilon
 =\frac{P_{t,T}^\varepsilon(e^{\theta\Phi})(X_t^\varepsilon)}
 {P_{0,T}^\varepsilon(e^{\theta\Phi})(X_0^\varepsilon)}.
\]

\begin{theorem}[Convergent resolved likelihood, Girsanov, and BSDE diagram]
\label{thm:r3-c2-girsanov}
The likelihood martingales converge jointly with the A4 rough tangent to
\[
 M_t=
 \frac{P_{t,T}(e^{\theta\Phi})(X_t)}
 {P_{0,T}(e^{\theta\Phi})(X_0)}
 =\mathcal E\left(
 \int_0^t\theta\sigma^TDu(s,X_s)dW_s
 \right),
\]
where
\[
 u(t,x)=\theta^{-1}\log P_{t,T}(e^{\theta\Phi})(x).
\]
Under the limiting tilted law,
\[
 dX_t=(b+\theta A Du(t,X_t))dt+\sigma dW_t^\theta,
\]
and \((u(t,X_t),\sigma^TDu(t,X_t))\) is the unique bounded solution of the
quadratic BSDE.  The finite likelihoods are exact mean-one martingales and are
uniformly integrable.
\end{theorem}

\begin{proof}
Apply Lemma~\ref{lem:r4-c2-kernel} to
\(F=e^{\theta\Phi}\).  The common complex source neighborhood gives an
\(L^{1+\eta}\) bound, hence uniform integrability and convergence of the
normalized likelihood martingales.  The limiting Gaussian kernel makes \(u\)
a bounded classical solution of
\[
 \partial_tu+b\cdot Du+\tfrac12A:D^2u+
 \tfrac\theta2Du^TADu=0.
\]
It\^o's formula gives the stochastic exponential, Girsanov drift, and BSDE;
the exponential transform gives uniqueness.  All conditional convergence was
proved through resolved transition kernels before invoking stochastic
calculus.
\end{proof}

\subsection{Linearized history pressure and compressed memory}

Fix one exposed prepared phase \(\nu_\Psi\).  Let
\(P_t^{\Psi,\rm D}\) be its normalized Doob transition semigroup on histories
and let \(L_\Psi^{\rm D}\) be the corresponding generator on
\(L^2(\nu_\Psi)\).  For the normalized nonlinear history semigroup
\[
 \mathcal E_t^\Psi F
 =\log P_t^\Psi(e^F)-\log P_t^\Psi\mathbf1,
\]
its first derivative at zero is the linear normalized transition:
\[
 D\mathcal E_t^\Psi(0)A=P_t^{\Psi,\rm D}A.
\]
Let \(P\) be the finite resolved orthogonal projection.

\begin{theorem}[Memory is the Laplace transform of the pressure tangent]
\label{thm:r3-c2-memory}
For resolved \(A,B\in PL^2(\nu_\Psi)\) and \(\Re z>0\),
\[
 \int_0^\infty e^{-zt}
 \left\langle B,D\mathcal E_t^\Psi(0)A\right\rangle_{L^2(\nu_\Psi)}dt
 =\left\langle
 B,P(z-L_\Psi^{\rm D})^{-1}PA
 \right\rangle.
\]
Consequently the compressed memory transform is
\[
 \widehat K_\Psi(z)
 =zP-PL_\Psi^{\rm D}P-
 [P(z-L_\Psi^{\rm D})^{-1}P]^{-1}.
\]
Linearizing the nonlinear history pressure and then forming the compressed
memory gives the same operator as forming A4's compressed-resolvent identity
under the prepared phase.  A further source derivative on either route is the
same Duhamel insertion.
\end{theorem}

\begin{proof}
Differentiate the normalized Feynman--Kac formula at \(F=0\).  Division by
\(P_t^\Psi\mathbf1\) is exactly the Doob normalization, so the derivative is
\(P_t^{\Psi,\rm D}A\).  Pair with \(B\) and integrate its semigroup in time to
obtain the compressed resolvent.  A4 defines memory by the displayed Schur
complement, proving the first commutation statement.  Differentiating the
resolvent identity gives
\[
 D(z-L_\Psi^{\rm D})^{-1}
 =(z-L_\Psi^{\rm D})^{-1}
 (DL_\Psi^{\rm D})
 (z-L_\Psi^{\rm D})^{-1},
\]
which is the Laplace transform of the same Duhamel insertion obtained by
source-differentiating the history semigroup.
\end{proof}

\begin{theorem}[Round-four C2 closure]
\label{thm:r3-c2-main}
The series has a platform-labelled contraction functor, verified mechanical
maps, a strict path-potential dual with countably additive phase laws, a closed
coboundary cotangent quotient, model-derived first and second pressure
responses, operational entropic calibration, a proved deterministic-to-
diffusion likelihood/BSDE limit, and an exact memory--pressure commutative
identity.
\end{theorem}

\begin{proof}
Combine Theorems~\ref{thm:r3-c2-functor},
\ref{thm:r3-c2-strict}, \ref{thm:r3-c2-cotangent},
\ref{thm:r3-c2-response}, \ref{thm:r3-c2-rigidity},
\ref{thm:r3-c2-girsanov}, and \ref{thm:r3-c2-memory}.
\end{proof}
